problem:  
Let $\{(M_i^n,g_i)\}$ be a sequence of closed Riemannian $n$‑manifolds with $\mathrm{sec}_{g_i}\ge -1$ and $\mathrm{diam}(M_i,g_i)\le1$, collapsing in the Gromov–Hausdorff sense to a compact Alexandrov space $(X^k,d_X)$ with $k<n$.  
For each $x\in X$, fix a sequence of points $p_i\in M_i$ converging to $x$. The metric tangent cones $T_xX$ are known to be metric cones over the spaces of directions $\Sigma_x$, each an Alexandrov space of dimension $k-1$ and curvature $\ge1$.  

**Question:**  
Is there a finite collection of topological types $\mathcal{T}(n)$ with the following universal property?  
Whenever $\{(M_i^n,g_i)\}$ satisfies the above assumptions, there exists a subsequence and possibly passing to finite covering spaces $\widehat M_i\to M_i$ such that each blow‑up limit (space of directions) $\Sigma_x$ of the Gromov–Hausdorff limit $X$ has a topology belonging to $\mathcal{T}(n)$.  

Equivalently: *Does collapsing with a uniform lower sectional curvature bound admit a universal finite topological model collection for all spaces of directions appearing in the metric limit?*  

Why is it a "good" problem:  
This refined problem avoids the dimensional triviality inherent in tangent cones and instead focuses on the *topological types* of their cross‑sections, which remain largely mysterious in lower‑curvature collapses. In known two‑sided curvature settings, local collapsing is modeled by nilpotent torus bundles or finite quotients thereof, but under one‑sided curvature bounds, Alexandrov limit spaces may exhibit wildly varying link topologies.  
Establishing whether there exists a finite topological model collection $\mathcal{T}(n)$ would expose deep structural patterns—or the lack thereof—in the geometry of Alexandrov limits, extending ideas from Cheeger–Gromov theory, Perelman’s stability, and stratified Alexandrov geometry. It would also connect collapse geometry to topological finiteness phenomena and singularity theory. The problem is simply stated, logically sound, and targets a fundamental gap between smooth collapsing theory and the topology of singular metric limits—making it a genuinely “good” research-level question.